\documentclass[11pt,a4paper]{article}
\usepackage[utf8]{inputenc}
\usepackage[T1]{fontenc}
\usepackage[catalan]{babel}
\usepackage{lmodern}
\linespread{1.22}
\usepackage[a4paper,top=2.4cm,bottom=2.2cm,left=1.5cm,right=1.5cm,headheight=1.1cm,headsep=0.5cm,footskip=1cm]{geometry}
\usepackage{amsmath,amssymb,cancel}
\usepackage{icomma}
\usepackage{graphicx}
\usepackage[table]{xcolor}
\usepackage{tikz}\usetikzlibrary{calc}
\usepackage{pgfplots}\pgfplotsset{compat=1.17}
\usepackage{fontawesome5}
\usepackage{ccicons}
\usepackage{fancyhdr}
\usepackage{enumitem}
\usepackage{booktabs}
\usepackage{tabularx}
\usepackage{needspace}
\usepackage[most]{tcolorbox}
\usepackage{titlesec}
\usepackage[hidelinks]{hyperref}
\hypersetup{pdftitle={Apunts de teoria. Unitat 2: Àlgebra. Equacions i inequacions. 1r de Batxillerat MI},pdfsubject={Matemàtiques I},pdfkeywords={polinomis, identitats notables, regla de Ruffini, teorema del residu, factorització, fraccions algebraiques, equacions, inequacions}}

% colors de la web (FentMates, bloc d'Àlgebra: blau; es manté el nom fmgreen dels apunts de les altres unitats)
\definecolor{fmnavy}{HTML}{1B2748}
\definecolor{fmgreen}{HTML}{3D6AA0}
\definecolor{fmgreenl}{HTML}{4E81BD}
\definecolor{fmyellow}{HTML}{FFC94D}
\definecolor{fmgold}{HTML}{E9A919}
\definecolor{fmaccent}{HTML}{FF7A50}
\definecolor{fmteal}{HTML}{2BB3A3}
\definecolor{fmdef}{HTML}{1D4F73}\definecolor{fmdefbg}{HTML}{EEF4F9}
\definecolor{fmprop}{HTML}{3B4252}\definecolor{fmpropbg}{HTML}{F3F4F7}
\definecolor{fmex}{HTML}{A3203A}\definecolor{fmexbg}{HTML}{FBF0F2}
\definecolor{fmmuted}{HTML}{5D6685}\definecolor{fmball0}{HTML}{E2483D}\definecolor{fmball1}{HTML}{2F6FD6}
\definecolor{fmsurface}{HTML}{EEF1F8}
\definecolor{fmA}{HTML}{3B82F6}\definecolor{fmB}{HTML}{F97316}\definecolor{fmR}{HTML}{E11D48}\definecolor{fmK}{HTML}{16A34A}\definecolor{fmV}{HTML}{7C3AED}

\newcommand{\lt}{<}\newcommand{\gt}{>}
% fórmula en línia que no es parteix mai; si no cap en una línia, es posa centrada
\newcommand{\fmm}[1]{\setbox0=\hbox{$#1$}\ifdim\wd0>\linewidth\fmdisp{#1}\else\mbox{\box0}\fi}
% fórmula centrada; si és més ampla que la línia, s'encongeix fins que hi cap
\newcommand{\fmdisp}[1]{\[\setbox2=\hbox{$\displaystyle #1$}\ifdim\wd2>\linewidth\resizebox{\linewidth}{!}{\box2}\else\box2\fi\]}
\DeclareUnicodeCharacter{2212}{-}
\DeclareUnicodeCharacter{2225}{\ensuremath{\parallel}}\DeclareUnicodeCharacter{22A5}{\ensuremath{\perp}}\DeclareUnicodeCharacter{2226}{\ensuremath{\nparallel}}
\DeclareUnicodeCharacter{207A}{\ensuremath{^{+}}}\DeclareUnicodeCharacter{207B}{\ensuremath{^{-}}}
\setlength{\parindent}{0pt}
\setlength{\parskip}{12pt}
\setlength{\lineskip}{5pt}\setlength{\lineskiplimit}{5pt}  % més aire al voltant de les fraccions en línia
\setlength{\emergencystretch}{3em}
\setlist{itemsep=4pt,topsep=4pt}
\renewcommand{\thesection}{2.\arabic{section}}
\titleformat{\section}{\color{fmnavy}\Large\bfseries}{\colorbox{fmdef}{\color{white}\strut\ \thesection\ }}{0.6em}{}[{\color{fmdef}\titlerule[1.2pt]}]
\titlespacing*{\section}{0pt}{10pt}{4pt}
\titleformat*{\subsection}{\color{fmgreen}\large\bfseries}
\titlespacing*{\subsection}{0pt}{6pt plus 2pt}{0pt}
\allowdisplaybreaks

% dibuix de la targeta de la unitat a la web (illustrations.js), en blanc
\newcommand{\fmcard}{\begin{tikzpicture}[x=0.062cm,y=-0.062cm,white,line width=3pt,line cap=round]
  \begin{scope}[transparency group,opacity=0.15]
  \node[anchor=base west,inner sep=0pt,font=\fontsize{26}{26}\selectfont] at (14,44) {$P(x)=x^2+3x-4$};
  \draw[line width=2.4pt] (40,62) rectangle (84,106) (88,62) rectangle (106,106);
  \fill[opacity=0.3] (40,62) rectangle (84,106);
  \node[anchor=base west,inner sep=0pt,font=\fontsize{17}{17}\selectfont] at (52,90) {$x^2$};
  \node[anchor=base west,inner sep=0pt,font=\fontsize{17}{17}\selectfont] at (91,90) {$x$};
  \draw[line width=3.2pt] (140,106) .. controls (160,106) and (166,66) .. (186,66) .. controls (206,66) and (214,100) .. (226,74);
  \draw[line width=2pt,opacity=0.6] (134,106)--(230,106);
  \end{scope}
\end{tikzpicture}}
\newcommand{\fmbrand}{\raisebox{-0.25\height}{\includegraphics[height=0.5cm]{logo-fentmates.pdf}}}

% capçalera i peu
\pagestyle{fancy}\fancyhf{}
\fancyhead[L]{\small\fmbrand}
\fancyhead[R]{\small\color{fmmuted}1r Batxillerat MI · Unitat 2: Àlgebra. Equacions i inequacions}
\fancyfoot[L]{\scriptsize\color{fmmuted}\href{https://creativecommons.org/licenses/by-nc-sa/4.0/deed.ca}{\ccbyncsa}\ \ FentMates · fentmates.org · CC BY-NC-SA 4.0}
\fancyfoot[R]{\small\color{fmmuted}\thepage}
\renewcommand{\headrulewidth}{0.4pt}\renewcommand{\headrule}{\color{fmmuted!40}\hrule height \headrulewidth}
\fancypagestyle{primera}{\fancyhf{}\renewcommand{\headrulewidth}{0pt}
  \fancyfoot[L]{\scriptsize\color{fmmuted}\href{https://creativecommons.org/licenses/by-nc-sa/4.0/deed.ca}{\ccbyncsa}\ \ FentMates · fentmates.org · CC BY-NC-SA 4.0}
  \fancyfoot[R]{\small\color{fmmuted}\thepage}}

% caixes (mateixos colors que la web)
% quan una caixa es parteix entre dues pàgines, la vora del tall és discontínua
\tcbset{fmbase/.style={enhanced,breakable,arc=3mm,boxrule=1pt,left=3.5mm,right=3.5mm,top=2mm,bottom=2.5mm,before skip=10pt,after skip=10pt,fonttitle=\bfseries,
  frame style={fill=tcbcolframe},
  overlay first={\fill[tcbcolback] ([xshift=1pt]frame.south west) rectangle ([xshift=-1pt,yshift=1.2pt]frame.south east);
    \draw[tcbcolframe,line width=1pt,dash pattern=on 5pt off 3pt] (frame.south west)--(frame.south east);},
  overlay middle={\fill[tcbcolback] ([xshift=1pt]frame.south west) rectangle ([xshift=-1pt,yshift=1.2pt]frame.south east);
    \fill[tcbcolback] ([xshift=1pt]frame.north west) rectangle ([xshift=-1pt,yshift=-1.2pt]frame.north east);
    \draw[tcbcolframe,line width=1pt,dash pattern=on 5pt off 3pt] (frame.south west)--(frame.south east) (frame.north west)--(frame.north east);},
  overlay last={\fill[tcbcolback] ([xshift=1pt]frame.north west) rectangle ([xshift=-1pt,yshift=-1.2pt]frame.north east);
    \draw[tcbcolframe,line width=1pt,dash pattern=on 5pt off 3pt] (frame.north west)--(frame.north east);}}}
\newtcolorbox{fmdef}[1]{fmbase,colback=fmdefbg,colframe=fmdef,coltitle=white,title={#1}}
\newtcolorbox{fmprop}[1]{fmbase,colback=fmpropbg,colframe=fmprop,coltitle=white,title={#1},halign title=center}
\newtcolorbox{fmex}[1]{fmbase,before upper=\raggedright,colback=fmexbg,colframe=fmex,colbacktitle=fmexbg,coltitle=fmex,title={\faPencil*\ #1},
  titlerule=0.4pt,titlerule style={fmex!40,dashed}}
\newtcolorbox{fmwarn}[1]{fmbase,colback=fmsf!6,colframe=fmsf,coltitle=white,title={#1}}
\newtcolorbox{fmnote}[1][]{enhanced,breakable,colback=white,colframe=white,borderline west={3pt}{0pt}{fmaccent},arc=0pt,left=3mm,right=2mm,top=1mm,bottom=1mm,before skip=6pt,after skip=6pt,before upper=\raggedright,#1}
\newcommand{\fmkey}[1]{\colorbox{fmyellow!45}{\strut #1}}
\newcommand{\fmuline}[1]{\textbf{\color{fmaccent}#1}}
\definecolor{fmok}{HTML}{15803D}\definecolor{fmev}{HTML}{2563EB}\definecolor{fmsf}{HTML}{C2410C}\definecolor{fmsi}{HTML}{7C3AED}\definecolor{fmas}{HTML}{A3203A}
\newcommand{\fmbadge}[2]{\tcbox[on line,colback=#1,colframe=#1,arc=2mm,boxsep=0pt,left=4pt,right=4pt,top=1.5pt,bottom=1.5pt]{\color{white}\footnotesize\bfseries #2}}
\newcommand{\fmchip}[1]{\tcbox[on line,colback=fmsf!10,colframe=fmsf!50,arc=2mm,boxsep=1pt,left=4pt,right=4pt]{#1}}
\newcommand{\fmlevel}[2]{\tcbox[on line,colback=white,colframe=fmmuted!40,arc=2mm,boxsep=2pt,halign=center]{\begin{tabular}{c}{\scriptsize\color{fmmuted}\MakeUppercase{#1}}\\#2\end{tabular}}}
% gràfiques
\pgfplotsset{fmaxis/.style={axis lines=middle,grid=major,grid style={fmmuted!15},axis line style={fmmuted,-latex},
  tick label style={font=\scriptsize,color=fmmuted},clip=true,unbounded coords=jump,restrict y to domain=-200:200,
  every axis plot/.append style={line join=round}},
  fmcurve/.style={fmaccent,line width=1.6pt,smooth}}
\tikzset{fmasy/.style={fmmuted,dashed,line width=0.9pt},fmguide/.style={fmaccent,dotted,line width=0.8pt},
  fmdot/.style={circle,fill=fmaccent,inner sep=0pt,minimum size=5pt},fmopen/.style={circle,draw=fmaccent,fill=white,line width=1pt,inner sep=0pt,minimum size=5pt}}
\newcommand{\fmpau}[1]{\tcbox[on line,colback=fmnavy,colframe=fmnavy,arc=2mm,boxsep=0pt,left=4pt,right=4pt,top=1.5pt,bottom=1.5pt]{\color{white}\scriptsize\bfseries\MakeUppercase{#1}}\ }

\begin{document}
\thispagestyle{primera}
\begin{tikzpicture}[remember picture,overlay]
  \fill[fmgreen] (current page.north west) rectangle ([yshift=-10.1cm]current page.north east);
  \begin{scope}
    \clip (current page.north west) rectangle ([yshift=-10.1cm]current page.north east);
    \node[anchor=north west,inner sep=0pt] at ([xshift=6.6cm,yshift=-1.3cm]current page.north west) {\fmcard};
  \end{scope}
  \node[anchor=north west,inner sep=0pt] at ([xshift=1.5cm,yshift=-1.1cm]current page.north west) {\tcbox[colback=white,colframe=white,arc=2.5mm,boxsep=0pt,left=7pt,right=7pt,top=5pt,bottom=5pt]{\includegraphics[height=0.75cm]{logo-fentmates.pdf}}};
\end{tikzpicture}
\vspace*{0.0cm}
\tcbox[colback=white!15!fmgreen,colframe=white!15!fmgreen,arc=2.5mm,boxsep=0pt,left=6pt,right=6pt,top=3pt,bottom=3pt]{\color{white}\small\bfseries Apunts de teoria · 1r Batxillerat · Matemàtiques I}

\vspace{4pt}
{\raggedright\color{white}\fontsize{25}{30}\selectfont\bfseries{\color{fmyellow}Unitat 2} · Àlgebra. Equacions i inequacions\par}
\vspace{6pt}
\begin{minipage}{\linewidth}\color{white!88!fmgreen}\small Teoria i exemples resolts pas a pas: polinomis i les seues operacions, la regla de Ruffini i el teorema del residu, la factorització, les fraccions algebraiques, les equacions (polinòmiques, racionals, amb radicals, exponencials i logarítmiques), les inequacions amb una incògnita i la resolució de problemes.\end{minipage}\par
% l'índex comença sempre a la mateixa altura, per davall de la franja (com en les unitats 5-10)
\vspace*{\dimexpr 8.6cm-\pagetotal\relax}

\begin{tcolorbox}[colback=fmsurface,colframe=fmsurface,arc=3mm,left=4mm,right=4mm,top=2mm,bottom=2mm]
\small\textbf{Índex}\par\smallskip
\begin{minipage}[t]{0.49\linewidth}\small \mbox{\hyperlink{sec1}{\textbf{\color{fmgreen}2.1}\ \ Polinomis. Operacions}}\\[3pt]
\mbox{\hyperlink{sec2}{\textbf{\color{fmgreen}2.2}\ \ Ruffini i teorema del residu}}\\[3pt]
\mbox{\hyperlink{sec3}{\textbf{\color{fmgreen}2.3}\ \ Factorització de polinomis}}\\[3pt]
\mbox{\hyperlink{sec4}{\textbf{\color{fmgreen}2.4}\ \ Fraccions algebraiques}}\\[3pt]
\mbox{\hyperlink{sec5}{\textbf{\color{fmgreen}2.5}\ \ Equacions polinòmiques}}\\[3pt]
\mbox{\hyperlink{sec6}{\textbf{\color{fmgreen}2.6}\ \ Equacions racionals}}\\[3pt]\end{minipage}\hfill\begin{minipage}[t]{0.49\linewidth}\small \mbox{\hyperlink{sec7}{\textbf{\color{fmgreen}2.7}\ \ Equacions amb radicals}}\\[3pt]
\mbox{\hyperlink{sec8}{\textbf{\color{fmgreen}2.8}\ \ Equacions exponencials}}\\[3pt]
\mbox{\hyperlink{sec9}{\textbf{\color{fmgreen}2.9}\ \ Equacions logarítmiques}}\\[3pt]
\mbox{\hyperlink{sec10}{\textbf{\color{fmgreen}2.10}\ \ Inequacions amb una incògnita}}\\[3pt]
\mbox{\hyperlink{sec11}{\textbf{\color{fmgreen}2.11}\ \ Resolució de problemes}}\\[3pt]\end{minipage}
\end{tcolorbox}

\Needspace{8\baselineskip}
\section{Polinomis. Operacions}\hypertarget{sec1}{}
\begin{fmdef}{Definicions}
\begin{itemize}
\item Un \textbf{monomi} és el producte d'un nombre (\textbf{coeficient}) i una o més lletres elevades a exponents naturals (\textbf{part literal}). El \textbf{grau} d'un monomi és la suma dels exponents de les lletres. Per exemple, \fmm{-5x^3} té coeficient \fmm{-5} i grau 3. Dos monomis són \textbf{semblants} si tenen la mateixa part literal; només es poden sumar o restar monomis semblants.
\item Un \textbf{polinomi} en \fmm{x} és una suma de monomis: \fmm{\ P(x)=a_nx^n+a_{n-1}x^{n-1}+\cdots+a_1x+a_0,\ \ a_n\neq0.} El seu \textbf{grau} és \fmm{n}; \fmm{a_n} és el \textbf{coeficient principal} i \fmm{a_0}, el \textbf{terme independent}. És \textbf{complet} si té tots els termes des del grau \fmm{n} fins al terme independent i està \textbf{ordenat} si els termes van de major a menor grau.
\item El \textbf{valor numèric} de \fmm{P(x)} per a \fmm{x=a} és el nombre \fmm{P(a)} que s'obté en substituir \fmm{x} per \fmm{a}.
\item Dos polinomis són \textbf{iguals} si tenen el mateix grau i els coeficients dels termes del mateix grau són iguals. Això permet calcular coeficients desconeguts \emph{igualant coeficients}.
\end{itemize}
\end{fmdef}

\Needspace{5\baselineskip}
\begin{fmprop}{Operacions amb polinomis}
\begin{itemize}
\item \textbf{Suma i resta}: se sumen (o resten) els monomis semblants.\newline \fmm{\ \text{grau}(P\pm Q)\le\max(\text{grau}\,P,\ \text{grau}\,Q)}.
\item \textbf{Producte}: es multiplica cada monomi del primer per cada monomi del segon (propietat distributiva) i es redueixen els termes semblants. \fmm{\ \text{grau}(P\cdot Q)=\text{grau}\,P+\text{grau}\,Q}.
\end{itemize}

\textbf{Identitats notables}

\begin{center}\renewcommand{\arraystretch}{1.4}\begin{tabularx}{\linewidth}{XX}
$\displaystyle (a+b)^2=a^2+2ab+b^2$ & $\displaystyle (a-b)^2=a^2-2ab+b^2$\\[4pt]
$\displaystyle (a+b)(a-b)=a^2-b^2$ & $\displaystyle (a+b+c)^2=a^2+b^2+c^2+2ab+2ac+2bc$\\[4pt]
$\displaystyle (a+b)^3=a^3+3a^2b+3ab^2+b^3$ & $\displaystyle (a-b)^3=a^3-3a^2b+3ab^2-b^3$\\[4pt]
\end{tabularx}\end{center}
\end{fmprop}

\Needspace{5\baselineskip}
\begin{fmex}{Exemple resolt: operacions amb polinomis}
\begin{enumerate}[label=\textbf{\arabic*.},leftmargin=*]
\item Simplifica \fmm{3(x-1)^2-(2x+3)(2x-3)+x(x^2-4)}.

Desenvolupem amb les identitats notables i la propietat distributiva: \fmm{\ 3(x^2-2x+1)-(4x^2-9)+x^3-4x}

\fmm{=3x^2-6x+3-4x^2+9+x^3-4x=x^3-x^2-10x+12}
\item Calcula \fmm{a} i \fmm{b} perquè \fmm{(x+a)(2x-3)=2x^2+bx-6}.

\fmm{(x+a)(2x-3)=2x^2+(2a-3)x-3a}

Igualem coeficients: \fmm{\ -3a=-6\Rightarrow a=2;\quad b=2a-3=1}
\end{enumerate}
\end{fmex}

\Needspace{7\baselineskip}
\subsection*{Divisió de polinomis}
\begin{fmdef}{Divisió entera de polinomis}
Donats dos polinomis \fmm{D(x)} (\textbf{dividend}) i \fmm{d(x)\neq0} (\textbf{divisor}), existeixen dos únics polinomis \fmm{c(x)} (\textbf{quocient}) i \fmm{r(x)} (\textbf{residu}) tals que

\fmdisp{D(x)=d(x)\cdot c(x)+r(x)\qquad \text{amb}\ \ \text{gr}\,r(x)\lt\text{gr}\,d(x)}

Si \fmm{r(x)=0}, la divisió és \textbf{exacta}: \fmm{d(x)} és un \textbf{divisor} de \fmm{D(x)}.

\begin{fmnote}\fmuline{Procediment}: s'ordenen els polinomis (deixant buit el lloc dels termes que falten), es divideix el primer terme del dividend entre el primer del divisor, es multiplica el resultat pel divisor i es resta. Es repeteix fins que el grau del que queda és menor que el del divisor.\end{fmnote}
\end{fmdef}

\Needspace{8\baselineskip}
\section{Regla de Ruffini. Teorema del residu}\hypertarget{sec2}{}
\begin{fmdef}{Regla de Ruffini}
Quan el divisor és de la forma \fmm{x-a}, la divisió es fa ràpidament amb la \textbf{regla de Ruffini}:

\begin{enumerate}
\item S'escriuen els coeficients del dividend, ordenats i \textbf{posant 0} en els termes que falten. A l'esquerra s'escriu el nombre \fmm{a}.
\item Es baixa el primer coeficient. Es multiplica per \fmm{a}, s'escriu davall del coeficient següent i se sumen. Es repeteix fins al final.
\item L'últim nombre és el \textbf{residu}; els anteriors són els coeficients del \textbf{quocient}, que té un grau menys que el dividend.
\end{enumerate}

\begin{fmnote}\fmuline{Atenció}: si el divisor és \fmm{x+3}, aleshores \fmm{a=-3}, perquè \fmm{x+3=x-(-3)}.\end{fmnote}
\end{fmdef}

\Needspace{5\baselineskip}
\begin{fmprop}{Divisor de la forma $ax+b$}
Com que \fmm{ax+b=a\left(x+\frac ba\right)}, fem Ruffini amb el nombre \fmm{-\frac ba}. El \textbf{residu és el mateix}, però el quocient obtingut s'ha de \textbf{dividir entre \fmm{a}}.

Per exemple, \fmm{(6x^3-x^2+4x-3):(2x-1)}: Ruffini amb \fmm{\frac12} dona \fmm{6x^2+2x+5} i residu \fmm{-\frac12}. Dividint entre 2: \fmm{\ c(x)=3x^2+x+\frac52,\ \ r=-\frac12}.
\end{fmprop}

\Needspace{5\baselineskip}
\begin{fmdef}{Teorema del residu i teorema del factor}
\textbf{Teorema del residu}: el residu de la divisió \fmm{P(x):(x-a)} és igual al valor numèric \fmm{P(a)}.

\fmdisp{r=P(a)}

\emph{Demostració}: \fmm{P(x)=(x-a)\,c(x)+r}. Si substituïm \fmm{x=a}: \fmm{P(a)=0\cdot c(a)+r=r}.

Un nombre \fmm{a} és una \textbf{arrel} (o zero) de \fmm{P(x)} si \fmm{P(a)=0}. \textbf{Teorema del factor}:

\fmdisp{P(a)=0\iff P(x)\ \text{és divisible per}\ (x-a)\iff P(x)=(x-a)\cdot c(x)}

\end{fmdef}

\Needspace{5\baselineskip}
\begin{fmex}{Exemple resolt: teorema del residu}
\begin{enumerate}[label=\textbf{\arabic*.},leftmargin=*]
\item Calcula, sense fer la divisió, el residu de \fmm{(x^5-3x^2+7):(x+1)}.

\fmm{r=P(-1)=(-1)^5-3(-1)^2+7=-1-3+7=3}
\item Troba \fmm{k} perquè \fmm{P(x)=x^3+kx^2-5x+6} siga divisible per \fmm{x-3}.

Ha de ser \fmm{P(3)=0}: \fmm{\ 27+9k-15+6=0}

\fmm{9k=-18\Rightarrow k=-2}
\end{enumerate}
\end{fmex}

\Needspace{8\baselineskip}
\section{Factorització de polinomis}\hypertarget{sec3}{}
\begin{fmdef}{Arrels i factorització}
\textbf{Factoritzar} un polinomi és escriure'l com a producte de polinomis del menor grau possible (\textbf{irreductibles}). Si \fmm{a_1,a_2,\ldots,a_n} són les arrels de \fmm{P(x)=a_nx^n+\cdots+a_0}, aleshores

\fmdisp{P(x)=a_n\,(x-a_1)(x-a_2)\cdots(x-a_n)}

\begin{itemize}
\item Un polinomi de grau \fmm{n} té, com a màxim, \fmm{n} arrels reals.
\item Si un factor \fmm{(x-a)} apareix \fmm{k} vegades, \fmm{a} és una \textbf{arrel de multiplicitat \fmm{k}} (doble, triple\dots{}).
\item Les \textbf{arrels enteres} d'un polinomi amb coeficients enters són \textbf{divisors del terme independent} \fmm{a_0}. Les arrels racionals \fmm{\frac pq} (fracció irreductible) compleixen que \fmm{p} divideix \fmm{a_0} i \fmm{q} divideix \fmm{a_n}.
\item Hi ha polinomis sense arrels reals, com \fmm{x^2+1} o \fmm{x^2+x+1} \fmm{(\Delta\lt0)}: són irreductibles.
\end{itemize}
\end{fmdef}

\Needspace{5\baselineskip}
\begin{fmprop}{Procediment per a factoritzar}
\begin{enumerate}
\item \textbf{Traure factor comú} si és possible: \fmm{\ 2x^4-6x^3=2x^3(x-3)}.
\item Reconéixer \textbf{identitats notables}: \fmm{\ x^2-9=(x+3)(x-3),\ \ 4x^2+4x+1=(2x+1)^2}.
\item Si és de \textbf{grau 2}, resoldre l'equació: \fmm{\ ax^2+bx+c=a(x-x_1)(x-x_2)}.
\item Si és \textbf{biquadrat} \fmm{(ax^4+bx^2+c)}, fer el canvi \fmm{t=x^2}.
\item Si és de grau major, aplicar \textbf{Ruffini} provant els divisors del terme independent (i, si cal, les fraccions \fmm{\frac pq}) fins a arribar a grau 2.
\end{enumerate}
\end{fmprop}

\Needspace{5\baselineskip}
\begin{fmex}{Exemple resolt: factorització amb Ruffini}
\begin{enumerate}[label=\textbf{\arabic*.},leftmargin=*]
\item Factoritza \fmm{P(x)=x^4-x^3-7x^2+x+6}.

Possibles arrels enteres: divisors de 6: \fmm{\ \pm1,\pm2,\pm3,\pm6}.

Ruffini amb \fmm{x=1} (residu 0) i, amb el quocient, \fmm{x=-1} (residu 0): \fmm{\ P(x)=(x-1)(x+1)(x^2-x-6)}

\fmm{x^2-x-6=0\Rightarrow x=\dfrac{1\pm5}{2}\Rightarrow x=3,\ x=-2}

\fmm{P(x)=(x-1)(x+1)(x-3)(x+2)}
\item Factoritza \fmm{Q(x)=2x^4+x^3-5x^2+2x}.

Factor comú: \fmm{\ Q(x)=x(2x^3+x^2-5x+2)}

Ruffini amb \fmm{x=1}: \fmm{\ 2x^3+x^2-5x+2=(x-1)(2x^2+3x-2)}

\fmm{2x^2+3x-2=0\Rightarrow x=\dfrac{-3\pm5}{4}\Rightarrow x=\frac12,\ x=-2}. Aleshores \fmm{2x^2+3x-2=2\left(x-\frac12\right)(x+2)=(2x-1)(x+2) \Rightarrow Q(x)=x\,(x-1)(2x-1)(x+2)}
\end{enumerate}
\end{fmex}

\Needspace{8\baselineskip}
\section{Fraccions algebraiques}\hypertarget{sec4}{}
\begin{fmdef}{Definicions}
\begin{itemize}
\item Una \textbf{fracció algebraica} és un quocient de dos polinomis, \fmm{\frac{P(x)}{Q(x)}}, amb \fmm{Q(x)\neq0}. Només té sentit per als valors de \fmm{x} que \textbf{no anu\lgem{}en el denominador}.
\item Dues fraccions \fmm{\frac PQ} i \fmm{\frac RS} són \textbf{equivalents} si \fmm{P\cdot S=Q\cdot R}.
\item \textbf{Simplificar} és dividir numerador i denominador per un mateix factor: cal \textbf{factoritzar} i eliminar els factors comuns.
\item El \textbf{m.c.m.} de diversos polinomis és el producte dels factors comuns i no comuns elevats al major exponent; el \textbf{m.c.d.}, el producte dels factors comuns elevats al menor exponent.
\end{itemize}

\begin{fmnote}\fmuline{Atenció}: només es poden simplificar \emph{factors}, mai \emph{sumands}: \fmm{\ \dfrac{x+3}{x}\neq3}.\end{fmnote}
\end{fmdef}

\Needspace{5\baselineskip}
\begin{fmprop}{Operacions amb fraccions algebraiques}
\begin{itemize}
\item \textbf{Suma i resta}: es redueixen a \textbf{comú denominador} (m.c.m. dels denominadors) i se sumen o resten els numeradors. \fmm{\ \dfrac PQ\pm\dfrac RS=\dfrac{P\cdot S\pm Q\cdot R}{Q\cdot S}}
\item \textbf{Producte}: \fmm{\ \dfrac PQ\cdot\dfrac RS=\dfrac{P\cdot R}{Q\cdot S}\qquad} \textbf{Quocient}: \fmm{\ \dfrac PQ:\dfrac RS=\dfrac{P\cdot S}{Q\cdot R}}
\item És convenient \textbf{factoritzar abans d'operar} i simplificar al final.
\end{itemize}
\end{fmprop}

\Needspace{5\baselineskip}
\begin{fmex}{Exemple resolt: operacions amb fraccions}
\begin{enumerate}[label=\textbf{\arabic*.},leftmargin=*]
\item Simplifica \fmm{\dfrac{x^3-4x}{x^3-4x^2+4x}}.

\fmm{\dfrac{x(x^2-4)}{x(x^2-4x+4)}=\dfrac{x(x+2)(x-2)}{x(x-2)^2}=\dfrac{x+2}{x-2}}
\item Calcula \fmm{\dfrac{x}{x-1}-\dfrac{2}{x+1}+\dfrac{2}{x^2-1}}.

Factoritzem els denominadors: \fmm{\ x^2-1=(x+1)(x-1)\Rightarrow} m.c.m. \fmm{=(x+1)(x-1)}

\fmm{\dfrac{x(x+1)-2(x-1)+2}{(x+1)(x-1)}=\dfrac{x^2+x-2x+2+2}{(x+1)(x-1)}=\dfrac{x^2-x+4}{x^2-1}}
\item Calcula i simplifica \fmm{\dfrac{x^2-9}{x^2+2x}\cdot\dfrac{x+2}{x-3}} i \fmm{\dfrac{2x-4}{x^2}:\dfrac{x^2-4}{x}}.

\fmm{\dfrac{(x+3)(x-3)}{x(x+2)}\cdot\dfrac{x+2}{x-3}=\dfrac{(x+3)(x-3)(x+2)}{x(x+2)(x-3)}=\dfrac{x+3}{x}}

\fmm{\dfrac{2(x-2)}{x^2}\cdot\dfrac{x}{(x+2)(x-2)}=\dfrac{2}{x(x+2)}}
\end{enumerate}
\end{fmex}

\Needspace{8\baselineskip}
\section{Equacions polinòmiques}\hypertarget{sec5}{}
\begin{fmdef}{Equacions de primer i segon grau}
\begin{itemize}
\item \textbf{Primer grau}: \fmm{ax+b=0\Rightarrow x=-\frac ba}. Si hi ha denominadors, es multiplica tot pel m.c.m.
\item \textbf{Segon grau}: \fmm{\ ax^2+bx+c=0\ \Rightarrow\ x=\dfrac{-b\pm\sqrt{b^2-4ac}}{2a}}
\item El \textbf{discriminant} \fmm{\Delta=b^2-4ac} indica el nombre de solucions: \fmm{\Delta\gt0}: dues solucions; \fmm{\Delta=0}: una solució doble; \fmm{\Delta\lt0}: cap solució real.
\item \textbf{Incompletes}: \fmm{ax^2+bx=0\Rightarrow x(ax+b)=0\Rightarrow x=0,\ x=-\frac ba}; \fmm{\ ax^2+c=0\Rightarrow x=\pm\sqrt{-\frac ca}}.
\item Si \fmm{x_1} i \fmm{x_2} són les solucions: \fmm{x_1+x_2=-\frac ba,\ \ x_1\cdot x_2=\frac ca}.
\end{itemize}
\end{fmdef}

\Needspace{5\baselineskip}
\begin{fmdef}{Equacions biquadrades i de grau superior}
\begin{itemize}
\item \textbf{Biquadrades}: \fmm{ax^4+bx^2+c=0}. Es fa el canvi \fmm{t=x^2} i s'obté \fmm{at^2+bt+c=0}. Per cada solució \fmm{t\ge0} s'obtenen \fmm{x=\pm\sqrt t} (les \fmm{t\lt0} no donen solucions reals).
\item El mateix mètode serveix per a \fmm{ax^{2n}+bx^n+c=0} amb \fmm{t=x^n} (per exemple, \fmm{x^6-9x^3+8=0} amb \fmm{t=x^3}).
\item \textbf{Grau superior}: es passa tot a un membre, es \textbf{factoritza} (factor comú, Ruffini\dots{}) i s'aplica que un producte és zero si algun factor és zero:
\end{itemize}

\fmdisp{A(x)\cdot B(x)=0\iff A(x)=0\ \ \text{o}\ \ B(x)=0}

\end{fmdef}

\Needspace{5\baselineskip}
\begin{fmex}{Exemple resolt: equacions polinòmiques}
\begin{enumerate}[label=\textbf{\arabic*.},leftmargin=*]
\item Resol \fmm{x^4-5x^2+4=0}.

\fmm{t=x^2:\ \ t^2-5t+4=0\Rightarrow t=\dfrac{5\pm3}{2}\Rightarrow t=4,\ t=1}

\fmm{x^2=4\Rightarrow x=\pm2;\quad x^2=1\Rightarrow x=\pm1.\quad x\in\{-2,-1,1,2\}}
\item Resol \fmm{x^3-2x^2-5x+6=0}.

Per Ruffini, \fmm{x=1} és arrel: \fmm{\ (x-1)(x^2-x-6)=0}

\fmm{(x-1)(x-3)(x+2)=0\Rightarrow x=1,\ x=3,\ x=-2}
\end{enumerate}
\end{fmex}

\Needspace{8\baselineskip}
\section{Equacions racionals}\hypertarget{sec6}{}
\begin{fmdef}{Equacions racionals}
Són equacions en què la incògnita apareix en algun \textbf{denominador}. Per a resoldre-les:

\begin{enumerate}
\item Es factoritzen els denominadors i es calcula el m.c.m.
\item Es multipliquen els dos membres pel m.c.m. per a eliminar els denominadors.
\item Es resol l'equació polinòmica resultant.
\item \textbf{Es comproven les solucions}: cal descartar les que anu\lgem{}en algun denominador.
\end{enumerate}
\end{fmdef}

\Needspace{5\baselineskip}
\begin{fmex}{Exemple resolt: equacions racionals}
\begin{enumerate}[label=\textbf{\arabic*.},leftmargin=*]
\item Resol \fmm{\dfrac{x}{x-1}+\dfrac{2}{x+1}=\dfrac{2}{x^2-1}}.

m.c.m. \fmm{=(x+1)(x-1)}. Multipliquem els dos membres: \fmm{\ x(x+1)+2(x-1)=2}

\fmm{x^2+3x-4=0\Rightarrow x=\dfrac{-3\pm5}{2}\Rightarrow x=1,\ x=-4}

Comprovació: \fmm{x=1} anu\lgem{}a els denominadors \fmm{x-1} i \fmm{x^2-1}: es descarta. Solució: \fmm{x=-4}.
\item Resol \fmm{\dfrac{x}{x-2}-\dfrac{2}{x-2}=3}.

\fmm{x-2=3(x-2)\Rightarrow x-2=3x-6\Rightarrow x=2}

Però \fmm{x=2} anu\lgem{}a el denominador: l'equació \textbf{no té solució}.
\end{enumerate}
\end{fmex}

\Needspace{8\baselineskip}
\section{Equacions amb radicals}\hypertarget{sec7}{}
\begin{fmdef}{Equacions irracionals}
Són equacions en què la incògnita apareix dins d'una arrel. Per a resoldre-les:

\begin{enumerate}
\item S'\textbf{aïlla un radical} en un dels membres.
\item S'\textbf{eleven els dos membres} a l'índex de l'arrel (al quadrat si és arrel quadrada).
\item Si encara queda algun radical, es repeteix el procés.
\item Es resol l'equació obtinguda i \textbf{es comproven totes les solucions} en l'equació inicial.
\end{enumerate}

\begin{fmnote}\fmuline{Per què cal comprovar?} En elevar al quadrat poden aparéixer \textbf{solucions estranyes}, perquè \fmm{A^2=B^2} també es compleix quan \fmm{A=-B}. Amb arrels d'índex imparell no n'apareixen, però sempre és bo comprovar.\end{fmnote}
\end{fmdef}

\Needspace{5\baselineskip}
\begin{fmex}{Exemple resolt: equacions amb radicals}
\begin{enumerate}[label=\textbf{\arabic*.},leftmargin=*]
\item Resol \fmm{x-\sqrt{x+1}=5}.

Aïllem el radical: \fmm{\ x-5=\sqrt{x+1}}. Elevem al quadrat: \fmm{\ x^2-10x+25=x+1}

\fmm{x^2-11x+24=0\Rightarrow x=8,\ x=3}

Comprovació: \fmm{x=8:\ 8-\sqrt9=5\ \checkmark;\quad x=3:\ 3-\sqrt4=1\neq5\ \times}. Solució: \fmm{x=8}
\item Resol \fmm{\sqrt{2x+3}-\sqrt{x+1}=1}.

\fmm{\sqrt{2x+3}=1+\sqrt{x+1}\Rightarrow 2x+3=1+2\sqrt{x+1}+x+1}

\fmm{x+1=2\sqrt{x+1}\Rightarrow x^2+2x+1=4x+4\Rightarrow x^2-2x-3=0\Rightarrow x=3,\ x=-1}

Comprovació: \fmm{x=3:\ 3-2=1\ \checkmark;\quad x=-1:\ 1-0=1\ \checkmark}. Totes dues són solució.
\end{enumerate}
\end{fmex}

\Needspace{8\baselineskip}
\section{Equacions exponencials}\hypertarget{sec8}{}
\begin{fmdef}{Equacions exponencials}
Són equacions en què la incògnita apareix en un \textbf{exponent}. Es basen en la propietat:

\fmdisp{a^M=a^N\iff M=N\qquad(a\gt0,\ a\neq1)}

\textbf{Recorda} les propietats de les potències:
\fmdisp{a^m\cdot a^n=a^{m+n},\ \ \dfrac{a^m}{a^n}=a^{m-n},\ \ (a^m)^n=a^{m\cdot n},\ \ a^{-n}=\dfrac1{a^n},\ \ \sqrt[n]{a^m}=a^{m/n},\ \ a^0=1.}
\end{fmdef}

\Needspace{5\baselineskip}
\begin{fmprop}{Tipus d'equacions exponencials i com resoldre-les}
\begin{enumerate}
\item \textbf{Els dos membres es poden escriure com a potències de la mateixa base}: s'igualen els exponents.
\item \textbf{No es poden reduir a la mateixa base}: es prenen logaritmes en els dos membres: \fmm{\ a^x=b\Rightarrow x=\dfrac{\log b}{\log a}}.
\item \textbf{Sumes de potències de la mateixa base} amb exponents que difereixen en un nombre: es descomponen (\fmm{2^{x+1}=2\cdot2^x,\ 2^{x-1}=\frac{2^x}{2}}) i es trau factor comú \fmm{a^x}.
\item \textbf{Apareixen \fmm{a^{2x}} i \fmm{a^x}} (per exemple \fmm{4^x} i \fmm{2^x}): es fa el \textbf{canvi de variable} \fmm{t=a^x} (aleshores \fmm{a^{2x}=t^2}) i s'obté una equació de segon grau. Com que \fmm{t=a^x\gt0}, les solucions \fmm{t\le0} no són vàlides.
\end{enumerate}
\end{fmprop}

\Needspace{5\baselineskip}
\begin{fmex}{Exemple resolt: un de cada tipus}
\begin{enumerate}[label=\textbf{\arabic*.},leftmargin=*]
\item \fmm{8^{x+1}=4^{x^2}}

\fmm{2^{3x+3}=2^{2x^2}\Rightarrow 2x^2-3x-3=0\Rightarrow x=\dfrac{3\pm\sqrt{33}}{4}}
\item \fmm{5^x=12}

\fmm{x\log5=\log12\Rightarrow x=\dfrac{\log12}{\log5}\approx 1{,}544}
\item \fmm{3^{x+1}+3^x+3^{x-1}=39}

\fmm{3^x\left(3+1+\tfrac13\right)=39\Rightarrow 3^x\cdot\dfrac{13}{3}=39}

\fmm{3^x=9\Rightarrow x=2}
\item \fmm{4^x-6\cdot2^x+8=0}

Com que \fmm{4^x=(2^2)^x=(2^x)^2}, fem \fmm{t=2^x}: \fmm{\ t^2-6t+8=0\Rightarrow t=4,\ t=2}

\fmm{2^x=4\Rightarrow x=2;\quad 2^x=2\Rightarrow x=1}
\end{enumerate}
\end{fmex}

\Needspace{8\baselineskip}
\section{Equacions logarítmiques}\hypertarget{sec9}{}
\begin{fmdef}{Logaritmes}
El \textbf{logaritme} en base \fmm{a} \fmm{(a\gt0,\ a\neq1)} d'un nombre \fmm{N\gt0} és l'exponent al qual cal elevar \fmm{a} per a obtindre \fmm{N}:

\fmdisp{\log_aN=x\iff a^x=N}

\par\noindent \begin{minipage}[t]{0.48\linewidth}\vspace{0pt}\raggedright
\fmdisp{\begin{array}{l}\log_a(M\cdot N)=\log_aM+\log_aN\\[2pt] \log_a\dfrac MN=\log_aM-\log_aN\\[4pt] \log_aM^k=k\log_aM\end{array}}
\end{minipage}\hfill\begin{minipage}[t]{0.48\linewidth}\vspace{0pt}\raggedright
\fmdisp{\begin{array}{l}\log_aa=1\qquad\log_a1=0\\[4pt] \log_aN=\dfrac{\log N}{\log a}=\dfrac{\ln N}{\ln a}\end{array}}
\end{minipage}\par\medskip

\begin{fmnote}\fmuline{Atenció}: només existeix el logaritme dels nombres \textbf{positius}.\end{fmnote}
\end{fmdef}

\Needspace{5\baselineskip}
\begin{fmprop}{Com resoldre equacions logarítmiques}
\begin{enumerate}
\item Fent servir les propietats, s'agrupen els logaritmes fins a arribar a \fmm{\log_aA=\log_aB}, i aleshores \fmm{A=B}. Els nombres també es poden escriure com a logaritmes: \fmm{1=\log10,\ 2=\log100}\dots{}
\item O bé s'arriba a \fmm{\log_aA=k} i s'aplica la definició: \fmm{A=a^k}.
\item \textbf{Es comproven les solucions}: tots els logaritmes de l'equació original han de tindre arguments positius.
\end{enumerate}
\end{fmprop}

\Needspace{5\baselineskip}
\begin{fmex}{Exemple resolt: equacions logarítmiques}
\begin{enumerate}[label=\textbf{\arabic*.},leftmargin=*]
\item Resol \fmm{\log x+\log(x-3)=1}.

\fmm{\log\big(x(x-3)\big)=\log10\Rightarrow x^2-3x-10=0\Rightarrow x=5,\ x=-2}

Comprovació: \fmm{x=-2} no és vàlida (\fmm{\log(-2)} no existeix). Solució: \fmm{x=5}.
\item Resol \fmm{2\log x-\log(x+6)=0}.

\fmm{\log x^2=\log(x+6)\Rightarrow x^2-x-6=0\Rightarrow x=3,\ x=-2}

Només és vàlida \fmm{x=3}.
\end{enumerate}
\end{fmex}

\Needspace{8\baselineskip}
\section{Inequacions amb una incògnita}\hypertarget{sec10}{}
\begin{fmdef}{Definicions}
\begin{itemize}
\item Una \textbf{inequació} és una desigualtat entre dues expressions algebraiques \fmm{(\lt,\gt,\le,\ge)}. Les seues solucions solen ser \textbf{intervals}: \fmm{\ (a,b)=\{a\lt x\lt b\},\ \ [a,+\infty)=\{x\ge a\}}\dots{}
\item Si sumem o restem el mateix nombre als dos membres, la desigualtat es manté. Si \textbf{multipliquem o dividim per un nombre negatiu, la desigualtat canvia de sentit}: \fmm{-2x\lt6\Rightarrow x\gt-3}.
\end{itemize}
\end{fmdef}

\Needspace{5\baselineskip}
\begin{fmprop}{Com resoldre inequacions}
\begin{itemize}
\item \textbf{Lineals}: s'opera com en una equació de primer grau, tenint en compte la regla del signe.
\item \textbf{Polinòmiques}: es passa tot a un membre, \fmm{P(x)\gt0} (o \fmm{\lt,\le,\ge}); es calculen les arrels i es factoritza; les arrels divideixen la recta en intervals i s'estudia el signe de \fmm{P(x)} en cada un (taula de signes o provant un valor); es trien els intervals que compleixen la inequació. Els extrems s'inclouen si la desigualtat és \fmm{\le} o \fmm{\ge}.
\item \textbf{Racionals} \fmm{\frac{P(x)}{Q(x)}\ge0}: igual, amb les arrels del numerador i del denominador. \textbf{Les arrels del denominador no s'inclouen mai.} No es pot multiplicar per \fmm{Q(x)} perquè no en coneixem el signe.
\end{itemize}
\end{fmprop}

\Needspace{5\baselineskip}
\begin{fmex}{Exemple resolt: inequacions}
\begin{enumerate}[label=\textbf{\arabic*.},leftmargin=*]
\item Resol \fmm{x^3-x^2-6x\ge0}.

Factoritzem: \fmm{\ x(x^2-x-6)=x(x-3)(x+2)}. Arrels: \fmm{-2,\ 0,\ 3}.

\begin{center}\renewcommand{\arraystretch}{1.5}\begin{tabular}{cccccccc}\toprule
\textbf{} & \textbf{$(-\infty,-2)$} & \textbf{$-2$} & \textbf{$(-2,0)$} & \textbf{$0$} & \textbf{$(0,3)$} & \textbf{$3$} & \textbf{$(3,+\infty)$}\\\midrule
$x+2$ & \textcolor{fmas}{$\boldsymbol{-}$} & 0 & \textcolor{fmok}{\textbf{+}} &  & \textcolor{fmok}{\textbf{+}} &  & \textcolor{fmok}{\textbf{+}}\\
$x$ & \textcolor{fmas}{$\boldsymbol{-}$} &  & \textcolor{fmas}{$\boldsymbol{-}$} & 0 & \textcolor{fmok}{\textbf{+}} &  & \textcolor{fmok}{\textbf{+}}\\
$x-3$ & \textcolor{fmas}{$\boldsymbol{-}$} &  & \textcolor{fmas}{$\boldsymbol{-}$} &  & \textcolor{fmas}{$\boldsymbol{-}$} & 0 & \textcolor{fmok}{\textbf{+}}\\
$P(x)$ & \textcolor{fmas}{$\boldsymbol{-}$} & 0 & \textcolor{fmok}{\textbf{+}} & 0 & \textcolor{fmas}{$\boldsymbol{-}$} & 0 & \textcolor{fmok}{\textbf{+}}\\
\bottomrule\end{tabular}\end{center}

\fmm{x\in[-2,\,0]\cup[3,+\infty)}
\item Resol \fmm{\dfrac{x-1}{x+3}\le0}.

Arrel del numerador: \fmm{x=1} (s'inclou). Arrel del denominador: \fmm{x=-3} (s'exclou).

Signes: \fmm{(-\infty,-3):+;\ \ (-3,1):-;\ \ (1,+\infty):+\Rightarrow x\in(-3,\,1]}
\end{enumerate}
\end{fmex}

\Needspace{8\baselineskip}
\section{Resolució de problemes}\hypertarget{sec11}{}
\begin{fmprop}{Pautes per a resoldre problemes amb equacions i inequacions}
\begin{enumerate}
\item \textbf{Llegir} l'enunciat amb atenció i identificar les dades i el que es demana.
\item \textbf{Anomenar les incògnites} (\fmm{x,y,z}\dots{}) i indicar què representa cada una (i en quines unitats).
\item \textbf{Plantejar} l'equació o la inequació traduint l'enunciat al llenguatge algebraic.
\item \textbf{Resoldre'l.}
\item \textbf{Comprovar i interpretar} les solucions: descartar les que no tenen sentit en el context (valors negatius, no enters\dots{}) i respondre amb una frase.
\end{enumerate}
\end{fmprop}

\Needspace{5\baselineskip}
\begin{fmdef}{Alguns plantejaments habituals}
\begin{itemize}
\item \textbf{Percentatges}: el \fmm{p\,\%} de \fmm{x} és \fmm{\frac p{100}x}. Augmentar un \fmm{p\,\%}: multiplicar per \fmm{1+\frac p{100}}.
\item \textbf{Mescles}: (quantitat total)\fmm{\times}(preu de la mescla) \fmm{=} suma de (quantitat \fmm{\times} preu) de cada component.
\item \textbf{Treballs simultanis} (aixetes, fonts, persones): si un element tarda \fmm{t} hores a fer un treball, en una hora en fa \fmm{\frac1t}. Si treballen junts, se sumen les parts: \fmm{\ \dfrac1{t_1}+\dfrac1{t_2}=\dfrac1T}.
\item \textbf{Geometria}: teorema de Pitàgores, perímetres, àrees i volums.
\item \textbf{Creixement exponencial}: \fmm{N(t)=N_0\cdot a^t} o \fmm{N_0\cdot e^{kt}}; es resolen amb logaritmes.
\end{itemize}
\end{fmdef}

\Needspace{5\baselineskip}
\begin{fmex}{Exemple resolt: problemes}
\begin{enumerate}[label=\textbf{\arabic*.},leftmargin=*]
\item Una aixeta omple un dipòsit en 3 hores i una altra en 6 hores. Quant tardaran juntes?

En una hora omplin \fmm{\frac13+\frac16=\frac12} del dipòsit.

\fmm{\dfrac1T=\dfrac12\Rightarrow T=2} hores
\item La base d'un rectangle fa 3 cm més que l'altura i l'àrea és de 40 cm$^2$. Quines en són les dimensions?

\fmm{x=} altura: \fmm{\ x(x+3)=40\Rightarrow x^2+3x-40=0\Rightarrow x=5,\ x=-8}

La solució negativa no té sentit: l'altura fa 5 cm i la base, 8 cm.
\end{enumerate}
\end{fmex}

  
\end{document}
